\documentclass{report}
\usepackage[utf8]{inputenc}     % for éô
\usepackage[english]{babel}     % for proper word breaking at line ends
\usepackage[a4paper, left=1.5in, right=1.5in, top=1.5in, bottom=1.5in]{geometry}
                                % for page size and margin settings
\usepackage{imakeidx}           % for creating a look-up index
\usepackage{natbib}
\bibliographystyle{humannat}
\usepackage{graphicx}           % for image widths

%\graphicspath{{img/}{../img/}}
\usepackage{subfiles}

\usepackage{tikz}
\usetikzlibrary{positioning}
\usepackage[european,siunitx]{circuitikz}
\usepackage{pgfplots}
\usepgfplotslibrary{dateplot}
\usepgfplotslibrary{fillbetween}
\usetikzlibrary{calc}
\usepackage{adjustbox}
\usepackage[labelfont=bf]{caption}
\usepackage{subcaption}
\usepackage{amsmath,amssymb}    % for better equations
\usepackage{amsthm}             % for better theorem styles
\usepackage{mathtools}          % for greek math symbol formatting
\usepackage{enumitem}           % for control of 'enumerate' numbering
\usepackage{listings}           % for control of 'itemize' spacing
%\usepackage[obeyDraft]{todonotes}
                                % for clear TODO notes
\usepackage{todonotes}          % for clear TODO notes
% \usepackage{multicol}         % for creating a 2-column document
\usepackage[hidelinks]{hyperref}% page numbers and '\ref's become clickable
% \usepackage{mathabx}          % astro
\usepackage{fourier-orns}       % leaf symbol
\usepackage{marvosym}           % lightning symbol
\usepackage{siunitx}            % for writing physical units like \si{\micro\ohm\per\meter}
\usepackage{textcomp}
\usepackage{eurosym}
\DeclareSIUnit\EUR{\text{\euro}}
%\usepackage[parfill]{parskip}
\usepackage{empheq}
\setlength{\fboxsep}{1em}
\usepackage[normalem]{ulem}     % for strikethrough text
\usepackage{booktabs}
\usepackage{tabularx}
\renewcommand{\arraystretch}{1.2}
\usepackage{bold-extra}
\usepackage{bm}
\usepackage{upgreek}
\usepackage{dsfont}

% Font choice

\usepackage{comment}


%\newcommand{\mat}[1]{\ensuremath{\boldsymbol{\mathrm{#1}}}}
\newcommand{\mat}[1]{\ensuremath{\bm{\mathrm{#1}}}}

%\begin{comment}
\usepackage[T1]{fontenc}
\usepackage[boldsans]{ccfonts}

\renewcommand{\bfseries}{\fontfamily{bch}\fontseries{bx}\selectfont}
% \DeclareMathAlphabet{\mathbf} {OT1}{\mathsf}{bx}{n}
\DeclareMathAlphabet{\mathbf} {OT1}{bch}{bx}{n}


%\pdfpkmode{dpdfezzz}
%\pdfpkresolution=1200


\usepackage[pagestyles, clearempty]{titlesec}%
\usepackage{fmtcount,etoolbox}

%\begin{comment}
\titleformat{\chapter}[display]{\center}
{%\titlerule\vspace{1.4em}
\Large\it\chaptername{ }\numberstring{chapter}}
%{1\baselineskip}
{0pt}
{\Huge}
[]

\titlespacing*{\chapter}
    {0pt}% left
    {1cm}% before
    {0pt}% after

\titleformat{\part}[display]
{\centering\normalfont\huge}
{\partname{ }\thepart}
{0pt}
{\vspace{1em}\Huge}
[]

\titlespacing*{\part}{0pt}{8cm}{0pt}


%\renewcommand{\mat}[1]{\ensuremath{\mathbf{#1}}}
%\newcommand{\matSigma}{\Sigma}
%\end{comment}

\makeatletter
\renewcommand{\mat}[1]{
\StrLeft{#1}{1}[\chrfirst]
\IfSubStr{ABCDEFGHIJKLMNOPQRSTUVWXYZabcdefghijklmnopqrstuvwxyz}{\chrfirst}
   {\mathbf{#1}}
   {\bm{#1}}
}
\makeatother


%% DOCUMENT
\def\thesistitle{Power Grid Failures}
\def\thesissubtitle{Theory of Fluctuating Renewables}
\def\thesisauthorfirst{Fons}
\def\thesisauthorsecond{van der Plas}
\def\thesissupervisorfirst{dr. Henk}
\def\thesissupervisorsecond{Don}
\def\thesissecondreaderfirst{prof. dr. Eric}
\def\thesissecondreadersecond{Cator}
\def\thesisdate{July 2019}

%% FOR PDF METADATA
\title{\thesistitle}
\author{\thesisauthorfirst\space\thesisauthorsecond}
\date{\thesisdate}


%% TODO PACKAGE
\newcommand{\towrite}[1]{\todo[inline,color=yellow!10]{#1}}
\newcommand{\tocite}[1]{\todo[color=green!10]{#1}}
\newcommand{\todelete}[1]{\sout{#1}}

%% INDEX
\makeindex

%% THEOREM STYLES
\renewcommand\qedsymbol{\leafNE}

\newtheorem{theorem}{Theorem}[section]
\newtheorem{corollary}{Corollary}[theorem]
\newtheorem{lemma}[theorem]{Lemma}
\newtheorem{proposition}[theorem]{Proposition}

\theoremstyle{definition}
\newtheorem{definition}[theorem]{Definition}
\newtheorem{examplex}[theorem]{Example}
\newenvironment{example}
  {\renewcommand{\qedsymbol}{\textthing}%
   \pushQED{\qed}\begin{examplex}}
  {\popQED\end{examplex}}

\theoremstyle{remark}
\newtheorem*{remark}{Remark}
\newtheorem*{intuition}{Intuitive argument}

%% MATH STYLES
%\newcommand{\mat}[1]{\ensuremath{\mathbf{#1}}}
\newcommand{\mel}[1]{\ensuremath{{\mathrm{#1}}}}
\newcommand{\phym}[1]{\ensuremath{\mathsf{#1}}}
\newcommand{\range}[1]{\ensuremath{[#1]}}
\newcommand{\norm}[1]{\left\lVert#1\right\rVert}
\newcommand{\abs}[1]{\left\lvert#1\right\rvert}
\newcommand{\defeq}[0]{:=}
\newcommand{\LPF}[0]{\ensuremath{\text{\textup{LPF}}}}
% \newcommand{\FRT}[0]{\ensuremath{\text{\textup{FRT}}}}
\newcommand{\FRT}[0]{\ensuremath{\text{\textup{K}}}}

%% MATH OPERATORS
\newcommand{\conj}[1]{\overline{#1}}
\newcommand{\pointwise}{\ast}
\DeclareMathOperator{\diag}{diag}
\DeclareMathOperator{\Arg}{Arg}
\DeclareMathOperator{\Ima}{Im}
\DeclareMathOperator{\linspan}{Span}
\DeclareMathOperator{\gr}{gr}
\DeclareMathOperator{\nullity}{null}
\DeclareMathOperator{\rank}{rank}
\DeclareMathOperator{\PROB}{\mathbb{P}}
\DeclareMathOperator{\EXP}{\mathbb{E}}
\DeclareMathOperator{\COVMAT}{\mathcal{C}}
\DeclareMathOperator{\gaussdistr}{\mathcal{N}}
\DeclareMathOperator*{\argmin}{argmin}
\DeclareMathOperator*{\argmax}{argmax}
\DeclareMathOperator{\sign}{sign}

%% MISC
\newcommand{\define}[1]{\emph{#1}\index{#1}}
\newcommand{\ie}{i.e.\ }
\newcommand{\eg}{e.g.\ }
\newcommand{\almosteverywhere}{a.e.\ }

\setlength{\parindent}{0pt}
\setlength{\parskip}{.5em}

% https://tex.stackexchange.com/a/63549
\makeatletter
\def\oversortoftilde#1{\mathop{\vbox{\m@th\ialign{##\crcr\noalign{\kern3\p@}%
      \sortoftildefill\crcr\noalign{\kern3\p@\nointerlineskip}%
      $\hfil\displaystyle{#1}\hfil$\crcr}}}\limits}

\def\sortoftildefill{$\m@th \setbox\z@\hbox{$\braceld$}%
  \braceld\leaders\vrule \@height\ht\z@ \@depth\z@\hfill\braceru$}

\makeatother

\DeclareRobustCommand\inlineres{\raisebox{-.3mm}{\begin{circuitikz}\ctikzset{bipoles/length=7mm}\draw[] (0,0) to[R] (1,0);\end{circuitikz}}}
\DeclareRobustCommand\inlineac{\raisebox{-1.4mm}{\begin{circuitikz}\ctikzset{bipoles/length=7mm}\draw[] (0,0) to[sV] (1,0);\end{circuitikz}}}


\newcommand{\inlinebar}[2]{\begin{tikzpicture}
    \begin{axis}[
        grid=major,
        width=40mm,
        height=20mm,
        scale=1,
        xmin=0, xmax=1.2, ymin=0, ymax=1,
        xtick = {0,1}, ytick = {0.,1},
        xmajorticks=false,
        ymajorticks=false]
     \addplot[fill=#2, fill opacity=.3, draw=none, samples=500]
       plot (\x, {sign(#1-\x) } );
   \end{axis}
\end{tikzpicture}}


\newcommand{\inlineprob}[3]{\begin{tikzpicture}
    \begin{axis}[
        grid=major,
        width=40mm,
        height=20mm,
        scale=1,
        xmin=0, xmax=1.2, ymin=0, ymax=1,
        xtick = {0,1}, ytick = {0.,1},
        %xticklabel={\pgfmathparse{\tick*100}\pgfmathprintnumber{\pgfmathresult}\%},
        xmajorticks=false,
        ymajorticks=false]
     \addplot[fill=#3, fill opacity=.3, samples=500, draw=none]
       plot (\x, {1/(1+exp((\x-#1)*#2))} ) \closedcycle;
   \end{axis}
\end{tikzpicture}}

\newcommand{\TikzCross}{\mathbin{\tikz [x=1.4ex,y=1.4ex,line width=.2ex, red] \draw (0,0) -- (1,1) (0,1) -- (1,0);}}%

\definecolor{darkgreen}{rgb}{.1, .8, .1}
\newcommand{\Checkmark}{\color{darkgreen}\checkmark}

%%%%%%%%%%%%%%%%%%%%%%%






\begin{document}
\pagenumbering{roman}

%\begin{multicols}{2}
\begin{titlepage}
	\thispagestyle{empty}
	\newcommand{\HRule}{\rule{\linewidth}{0.5mm}}
	\center
	\textsc{\Large Radboud University Nijmegen}\\[.7cm]
	\includegraphics[width=25mm]{img/in_dei_nomine_feliciter.eps}\\[.5cm]
	\textsc{Faculty of Science}\\[0.5cm]
	
	\HRule \\[0.4cm]
	{ \huge \bfseries \thesistitle}\\[0.1cm]
	\textsc{\thesissubtitle}\\
	\HRule \\[.5cm]
	\textsc{\large Thesis BSc Mathematics}\\[.5cm]
	
	\begin{minipage}{0.4\textwidth}
	\begin{flushleft} \large
	\emph{Author:}\\
	\thesisauthorfirst\space \textsc{\thesisauthorsecond}
	\end{flushleft}
	\end{minipage}
	~
	\begin{minipage}{0.4\textwidth}
	\begin{flushright} \large
	\emph{Supervisor:} \\
	\thesissupervisorfirst\space \textsc{\thesissupervisorsecond} \\[1em]
	\emph{Second reader:} \\
	\thesissecondreaderfirst\space \textsc{\thesissecondreadersecond}
	\end{flushright}
	\end{minipage}\\[4cm]
	\vfill
	{\large \thesisdate}\\
	\clearpage
\end{titlepage}

\chapter*{Abstract}
\vspace{1cm}
Limiting global warming requires a rapid transition towards renewable generation of electric energy, creating new challenges for the security of supply. As global demand for electric power is increasing, transmission networks are being pushed to their limits. 
Sudden fluctuations in generation and load change the current through lines of the high-voltage transmission network, which might cause them to overload. In some cases, the failure of one line results in a cascade of other failures, eventually causing a blackout.

In this thesis, we will try to better understand the effect of fluctuations in solar and wind generation on the robustness of electric power transmission. To do so, we develop the mathematical theory used to model the transmission of stochastic generation.

Additionally, we show how this theory was used in a recent study \citep{Nesti2018emergentfailures} to identify the most vulnerable lines in the German transmission network using historical weather data for 2011, demonstrating both the usefulness and the limitations of our theory of fluctuating renewables.

\towrite{

By studying the effect of fluctuating renewables, we can better estimate which transmission lines are vulnerable to these emergent failures. The theory discussed in this thesis can be used to 

The aim of this bachelor thesis is to introduce the reader to the fascinating mathematics that are used to model the effect of \emph{fluctuating renewables}. 

Emergent failures can be prevented 

The transmission network is 

At the end of Part I, we extend our model to \emph{correlated, normally distributed generation}. 

In Part II, we demonstrate how the theory developed in Part I can be used to 

These procedures are applied to the SciGRID network: a digital representation of the German high-voltage transmission network (489 buses, 695 lines). We use the work of \cite{PyPSA} and  to estimate hourly load and generation series during the year 2011 for each bus. 

Following \cite{Nesti2018emergentfailures}, we estimate the covariances in (renewable) generation between buses, and we use linearised power flow physics to find a \emph{multivariate Gaussian distribution of line flows}. This allows us to identify lines with the highest probability of overloading.

Additionally, for each vulnerable line we can compute the \emph{most likely fluctuation} that would cause its emergent failure. 
When a transmission line fails, power flow is instantly redistributed among other lines in the network, which might cause them to also overload.
By estimating the most likely fluctuation, we can simulate these \emph{cascading failures} more realistically, to assess the damage that an emergent failure would cause.

The model that we 

To 

In Part I, we cover the mathematical theory (graph theory and probability theory) 

This bachelor thesis was written as a }

\chapter*{Notation}
\begin{align*}
\mat{Y}&\text{ is a matrix or vector} \\
\mel{Y}_{i,j}& \text{ is an element of a matrix or vector} \\
Y&\text{ is a (real or complex) number}\\
\phym{Y}&\text{ is a physical quantity (like `admittance')}\\
\conj{\mat{Y}}&\text{ is the \emph{element-wise complex conjugate} of $\mat{Y}$} \\
\mat{Y}^*&\text{ is the \emph{adjoint} of $\mat{Y}$ (which is equivalent to}\\
        &\quad \text{ the \emph{transpose} for real-valued $\mat{Y}$)} \\
\range{n}&=\{1,\dots,n\} \\
\text{\textthing}&\text{\emph{ end of example}} \\
\text{\leafNE}&\text{\emph{ end of proof}}
\end{align*}
\towrite{Een geel doosje geeft aan dat een stuk tekst \emph{nog geschreven moet worden}.}
\todo{Een oranje doosje is meer \emph{voor mezelf}, dingen waar ik later nog naar moet kijken.}
\tocite{Een groen doosje geeft \emph{bronnen} aan die relevant (kunnen) zijn.}

\setcounter{tocdepth}{1}
\tableofcontents

\clearpage
\pagenumbering{arabic}

%%%%%%%%%%%%%%%%%%%%%%%%%%%%%%%%%%%%%%%%%%%%
\part{Theory of Electric Power Networks}\label{part:theory}


\chapter{Graph theory}\label{chap:graphtheory}
\begin{circuitikz}[scale=\textwidth/15cm]\ctikzset{bipoles/length=7mm}\draw[] (0,0) to[R] (15,0);\end{circuitikz}
\vspace{1em}
\subfile{graph}

\chapter*{Circuit Theory}\label{chap:circuittheory}
\begin{circuitikz}[scale=\textwidth/15cm]\ctikzset{bipoles/length=7mm}\draw[] (0,0) to (15,0);\end{circuitikz}
\vspace{1em}
\subfile{circuit}

\chapter{Probability Theory}\label{chap:probtheory}
\begin{circuitikz}[scale=\textwidth/15cm]\ctikzset{bipoles/length=7mm}\draw[] (0,0) to[C] (15,0);\end{circuitikz}
\vspace{1em}
\subfile{probability}

\chapter{The Electric Grid}\label{chap:grid}
\begin{circuitikz}[scale=\textwidth/15cm]\ctikzset{bipoles/length=7mm}\draw[] (0,0) to[sV] (15,0);\end{circuitikz}
\vspace{1em}
\subfile{grid}

\chapter{Modelling Transmission Networks}\label{chap:model}
\begin{circuitikz}[scale=\textwidth/15cm]\ctikzset{bipoles/length=7mm}\draw[] (0,0) to[cute inductor] (15,0);\end{circuitikz}
\vspace{1em}
\subfile{model}

%%%%%%%%%%%%%%%%%%%%%%%%%%%%%%%%%%%%%%%%%%%%
\part{Predicting failures in Germany's grid}\label{part:sim}

\chapter{Methods}\label{chap:methods}
%\begin{circuitikz}[scale=\textwidth/15cm]\ctikzset{bipoles/length=7mm}
%\draw[] (7.5,0) node[npn, rotate=90] (thenpn) {};
%\draw[] (0,0) to (thenpn.C);
%\draw[] (thenpn.E) to (15,0);
%\end{circuitikz}
\begin{circuitikz}[scale=\textwidth/15cm]\ctikzset{bipoles/length=7mm}
\ctikzset{bipoles/cuteswitch/thickness=0.5}
\draw[] (0,0) to[cute open switch] (15,0);
\vspace{1em}
\end{circuitikz}

\subfile{methods}

\chapter{Results}\label{chap:results}
\begin{circuitikz}[scale=\textwidth/15cm]\ctikzset{bipoles/length=7mm}\draw[] (0,0) to[fuse] (15,0);\end{circuitikz}
\subfile{results}

\appendix
\chapter{Software}
\emph{This thesis would not have been possible without the (mostly open source) projects listed in this section. }
\begin{itemize}
\item 
\end{itemize}

Contributions:
\begin{itemize}
\item Together with Fabian Hofmann and the original authors, I have modified PyPSA's plotting method to allow for a custom geographical projection. (Transverse Mercator with centre )
\item A minimal version of this document . The title page is based on the work of Tom Salet.
\item 
\end{itemize}

%\section{hmmm}
%adf
%%\chapter{Aoasdf}
%asdf
%\chapter{verder}
%\section{tweede}
%%\chapter{Aoasdf}
%asdf

\bibliography{librarynew}
%\printindex
\end{document}
